\chapter{Introducción}\label{chap:introduccion}

La anemia falciforme constituye uno de los desafíos más significativos en salud pública global, representando una hemoglobinopatía hereditaria de alta carga sanitaria que compromete fundamentalmente el transporte de oxígeno en el organismo. Esta condición patológica ocasiona crisis vaso-oclusivas recurrentes que deterioran progresivamente la calidad de vida del paciente y reduce considerablemente la expectativa de vida de la población afectada. Las estimaciones epidemiológicas globales señalan más de 300.000 nacimientos anuales con la enfermedad, manifestando una prevalencia particularmente sostenida en regiones tropicales y subtropicales donde persiste la ventaja heterocigota frente a la malaria, un fenómeno evolutivo que ha mantenido la frecuencia del alelo mutante en estas poblaciones \cite{Arishi2021}. Esta distribución geográfica específica coincide precisamente con las regiones que presentan mayores limitaciones en infraestructura sanitaria, creando una paradoja donde la mayor incidencia de la enfermedad se concentra en las áreas con menor capacidad diagnóstica y terapéutica.

Las estrategias diagnósticas vigentes para la detección de anemia falciforme, que abarcan desde los tradicionales frotis sanguíneos teñidos hasta técnicas más sofisticadas como la cromatografía líquida de alta resolución (HPLC) y la electroforesis de hemoglobina, ofrecen niveles de sensibilidad y especificidad elevados que pueden superar el {95} en condiciones óptimas. Sin embargo, estas metodologías demandan laboratorios especializados con equipamiento costoso, personal altamente entrenado en técnicas de laboratorio clínico, y cadenas de suministro estables para reactivos y consumibles especializados, elementos que raramente convergen de manera sostenible en entornos rurales o con recursos limitados \cite{McGann2017,Shrestha2025}. Esta dependencia crítica de infraestructura compleja no solo retrasa significativamente la implementación de programas de cribado neonatal universal, sino que también compromete la continuidad del tratamiento y el seguimiento longitudinal de pacientes en países de ingresos bajos y medios, donde se concentra aproximadamente el 75% de los nacimientos con anemia falciforme a nivel mundial. La situación se agrava considerablemente cuando se considera que el diagnóstico temprano y la intervención terapéutica oportuna pueden reducir la mortalidad infantil asociada a la enfermedad en más del 70%, evidenciando la urgente necesidad de desarrollar alternativas diagnósticas más accesibles y adaptadas a contextos de recursos limitados.

En este contexto de necesidad crítica, la microscopía digital holográfica sin lente (DLHM, por sus siglas en inglés) surge como una plataforma óptica innovadora y asequible que revoluciona el paradigma tradicional de la microscopía. Esta tecnología emplea fuentes de luz coherente económicas, como diodos láser o LEDs parcialmente coherentes, combinadas con sensores de imagen comerciales ampliamente disponibles para registrar patrones de interferencia complejos en una sola captura, eliminando completamente la necesidad de lentes objetivos costosos y sistemas de alineación precisos \cite{Amann2019,BuitragoDuque2023}. Al operar en el régimen de difracción de Fresnel, el montaje experimental preserva tanto la información de fase como de amplitud en el holograma digital registrado, y habilita reconstrucciones numéricas posteriores de alta calidad cuando se dispone de los parámetros geométricos del sistema debidamente caracterizados \cite{Greenbaum2014,Kim2025}. No obstante, la reconstrucción convencional de estos hologramas digitales exige la implementación de transformadas de Fourier computacionalmente costosas, algoritmos iterativos de propagación de campo, y una caracterización meticulosa de distancias de propagación y dimensiones de aperturas que puede resultar técnicamente desafiante en entornos no especializados. La propuesta innovadora de clasificar directamente los hologramas sin necesidad de reconstrucción evita dichas restricciones técnicas y computacionales, reduce drásticamente el tiempo de procesamiento de minutos a segundos, y conserva íntegra la rica información interferométrica relevante para describir las morfologías anómalas características de los eritrocitos falciformes.

El vertiginoso auge de las técnicas de aprendizaje automático en la última década ha impulsado el desarrollo de soluciones computacionales sofisticadas para hematología digital que abarcan un amplio espectro tecnológico, desde arquitecturas neuronales profundas de gran escala con millones de parámetros hasta esquemas híbridos que combinan extracción manual de descriptores morfológicos con clasificadores de última generación \cite{OConnor2020,Sazak2024,Alzubaidi2020}. Aunque estos modelos de vanguardia alcanzan métricas de rendimiento notables que pueden superar el 98 de precisión en conjuntos de validación controlados, su adopción clínica práctica se ve severamente limitada por la necesidad de grandes volúmenes de datos anotados para entrenamiento, la opacidad inherente de sus procesos decisionales que dificulta la interpretación médica, y los requerimientos computacionales que exceden las capacidades disponibles en muchos entornos clínicos. Revisiones sistemáticas recientes enfatizan la pertinencia crítica de metodologías interpretables que puedan ser validadas rigurosamente en muestras pequeñas mediante protocolos estadísticos robustos, así como el reporte sistemático y transparente de métricas de incertidumbre y calibración que permitan evaluar la confiabilidad de las predicciones \cite{Vabalas2019,Tsamardinos2022,Montero2024}. En consecuencia, resulta estratégico desde una perspectiva de implementación clínica priorizar descriptores físicos comprensibles basados en principios ópticos establecidos y clasificadores clásicos bien caracterizados que mantengan un equilibrio óptimo entre desempeño diagnóstico, trazabilidad de decisiones, y coste computacional accesible.

El presente proyecto de investigación se apoya en un repositorio propio meticulosamente curado de 304 hologramas digitales rotulados por expertos hematólogos, compuesto por 151 muestras de individuos sanos y 153 muestras confirmadas con anemia falciforme mediante técnicas gold standard. Este conjunto de datos principal se complementa estratégicamente con 30 hologramas externos de organismos no sanguíneos, incluyendo microalgas, polen y otros especímenes biológicos, que representan escenarios fuera de dominio cruciales para evaluar la robustez del sistema. Sobre esta base experimental se diseña un pipeline modular altamente configurable que unifica los diversos orígenes de datos mediante protocolos estandarizados, normaliza la geometría de los hologramas para garantizar consistencia en el procesamiento, y extrae sistemáticamente 84 características cuantitativas de textura, frecuencia y morfología inspiradas en fenómenos físicos bien establecidos en la literatura óptica. El sistema implementado entrena un ensamble sofisticado de clasificadores lineales y basados en árboles de decisión, integra esquemas de validación cruzada estratificada y Leave-One-Out para gestionar eficientemente la limitada disponibilidad de muestras clínicas, e incorpora un detector de anomalías robusto soportado en la distancia de Mahalanobis para identificar muestras atípicas o fuera de distribución \cite{Mahalanobis1936,Rousseeuw2019}. Todo el conjunto de herramientas se encapsula en una interfaz de línea de comandos reproducible y documentada que automatiza el flujo completo del modelo final, genera reportes estandarizados en formatos múltiples, y facilita la auditoría exhaustiva de resultados por parte de personal clínico y técnico.

El tratamiento completamente digital de hologramas fomenta significativamente la colaboración científica internacional al permitir que los modelos entrenados, scripts de procesamiento, y conjuntos de datos derivados puedan compartirse con facilidad a través de repositorios en línea y plataformas de colaboración. La documentación exhaustiva de dependencias de software, procedimientos detallados de ejecución, y bitácoras completas de experimentos incrementa sustancialmente la capacidad de replicación independiente y acelera la validación cruzada por parte de otras instituciones académicas y centros de investigación \cite{OConnor2020}. El repositorio asociado a este trabajo consolida de manera organizada el código fuente completo, archivos de configuración parametrizables, y guías operativas bilingües en español e inglés, en estricta consonancia con las directrices de propiedad intelectual y difusión del conocimiento establecidas por la Universidad EAFIT \cite{EAFIT2018}. Este compromiso institucional con los principios de ciencia abierta busca deliberadamente allanar el camino hacia ensayos clínicos multicéntricos de mayor escala y facilitar la transferencia tecnológica efectiva hacia iniciativas concretas de telemedicina en contextos de bajos recursos, donde el impacto potencial de estas tecnologías es máximo.

Los aportes sustanciales del presente estudio se concentran en cuatro frentes metodológicos articulados entre sí que constituyen una contribución integral al campo. En primer lugar, se expone un marco teórico comprehensivo que relaciona sistemáticamente los fundamentos ópticos de DLHM con los requisitos clínicos específicos del diagnóstico de anemia falciforme, resaltando las implicaciones técnicas y prácticas de trabajar directamente con hologramas sin reconstrucción numérica. En segundo lugar, se documenta exhaustivamente un pipeline reproducible de procesamiento que automatiza la extracción de características discriminantes, implementa esquemas de validación estadísticamente rigurosos, y genera visualizaciones interpretables del modelo final que facilitan la comprensión clínica. En tercer lugar, se demuestra experimentalmente mediante análisis cuantitativos que un ensamble cuidadosamente diseñado de clasificadores clásicos puede lograr métricas competitivas con una fracción mínima del costo computacional asociado a alternativas de mayor complejidad arquitectónica, sin sacrificar la interpretabilidad esencial para la adopción clínica. Finalmente, se implementa un detector de anomalías innovador que robustece significativamente la toma de decisiones al señalar proactivamente entradas fuera de dominio y justificar la alerta mediante contribuciones cuantificables por característica, proporcionando así un mecanismo de seguridad crítico para aplicaciones clínicas.

El manuscrito se organiza sistemáticamente de la siguiente manera para presentar de forma coherente el desarrollo completo del proyecto. El Capítulo~\ref{chap:problema} delimita con precisión la pregunta de investigación central, establece los objetivos específicos y generales, y define claramente el alcance y las limitaciones del estudio. El Capítulo~\ref{chap:marco} desarrolla exhaustivamente el sustento teórico necesario, cubriendo tanto los principios fundamentales de óptica holográfica como los conceptos esenciales de inteligencia artificial aplicada al diagnóstico médico. El Capítulo~\ref{chap:metodologia} describe detalladamente la metodología experimental empleada para construir, entrenar y validar el sistema propuesto, incluyendo protocolos de adquisición de datos y esquemas de evaluación. El Capítulo~\ref{chap:desarrollo} resume las decisiones críticas de ingeniería de software, documenta la arquitectura modular del repositorio, y presenta las consideraciones de implementación práctica. El Capítulo~\ref{chap:resultados} analiza cuantitativamente los hallazgos experimentales obtenidos, evalúa el rendimiento del detector de anomalías, y discute las implicaciones clínicas de los resultados. Por último, el Capítulo~\ref{chap:conclusiones} presenta las conclusiones generales del estudio, identifica las limitaciones actuales del enfoque propuesto, y delinea las líneas prometedoras de trabajo futuro que podrían expandir el impacto de esta investigación.